\chapter{Limits and Continuity}
In the real world, you may know a ``limit'' in its English definition as a point or boundary which cannot be exceeded. In fact, in mathematics, limits are central to the definitions of what will make up the core functions of calculus.

For a mathematical function, we may take its ``limit'' by observing the function's output leading up to a certain point.

\section{The Limit of a Function}
Let us observe a function $f(x)=(x^2-4)/(x-2)$, graphed in Figure \ref{fig:function-example}. At the value $x=2$, the function is undefined, or $0/0$. This is represented by a hollow circle in the graph. Observe what happens as we \textit{approach} the value $x=2$.

\begin{center}
	\begin{tabular}{| >{\color{mathGreen}}c| >{\color{mathGreen}}c| >{\color{mathGreen}}c| >{\color{mathGreen}}c| >{\color{mathGreen}}c| >{\color{mathGreen}}c| >{\color{mathGreen}}c| >{\color{mathGreen}}c| >{\color{mathGreen}}c| >{\color{mathGreen}}c|}
		\hline
		$x$    & 1 & 1.5 & 1.9 & 1.99 & $\rightarrow 2 \leftarrow$ & 2.01 & 2.1 & 2.5 & 3 \\
		\hline
		$f(x)$ & 3 & 3.5 & 3.9 & 3.99 & ? & 4.01 & 4.1 & 4.5 & 5 \\
		\hline
	\end{tabular}
\end{center}

\begin{marginfigure}
	\centering
	\resizebox{\linewidth}{!}{%
		\begin{tikzpicture}[>=stealth, scale=0.8]
	% Axes
	\draw[->] (-3, 0) -- (4.5, 0) node[below] {$x$};
	\draw[->] (0, -1) -- (0, 6.5) node[left] {$y$};
	
	\draw[dashed, gray!70] (2,0) -- (2,4);
	\draw[dashed, gray!70] (0,4) -- (2,4);
	
	% Axis tick labels
	\draw (2, 0.1) -- (2, -0.1) node[below, font=\footnotesize] {$2$};
	\draw (0.1, 4) -- (-0.1, 4) node[left, font=\footnotesize] {$4$};
	
	% Graph line and discontinuity
	\draw[line width=1.2pt, mathPurple] (-3, -1) -- (1.92, 3.92);
	\draw[line width=1.2pt, mathPurple] (2.08, 4.08) -- (4.5, 6.5);
	\draw[line width=1.2pt, mathPurple, fill=white] (2,4) circle (2.5pt);
\end{tikzpicture}%
	}
	\caption{The graph of $f(x)$, where the function is undefined at $x=2$.}
	\label{fig:function-example}
\end{marginfigure}

It is clear that as $x$ \textit{approaches} 2, $f(x)$ \textit{approaches} 4. In essence, this is what the ``limit of a function'' means. When we are taking limits, we look at the behavior of a function as it \textit{approaches} a point, not the value at the point itself.

\begin{defn}[Limit of a Function]
	Let $f$ be a function defined on all $x$ on the open interval $I$ containing $a$, except possibly at $a$ itself. The limit of $f(x)$ as $x$ approaches $a$ is $L$, defined as
	\begin{equation}
		\lim_{x\to a} f(x) = L.
	\end{equation}
\end{defn}

Thus, for the example above we may simply state that the limit of the function $f(x) = (x^2-4)/(x-2)$ as $x$ approaches 2 is 4, or
\begin{equation}
	\lim_{x\to 2} \frac{x^2-4}{x-2} = 4.
\end{equation}

While this definition may seem grounded enough (it is, for the sake of an introduction), we can more rigorously define the limit using the \textit{Epsilon-Delta ($\epsilon-\delta$) definition of the limit} in the next section.

\section{Epsilon-Delta Definition of the Limit}
Unfortunately, the word \textit{approaches} is not good enough for rigor in mathematical description. What does approaching even mean? \textit{How} does $x$ approach $a$?

By rephrasing the earlier definition, we may say that \textit{if $x$ is within a certain tolerance level of $a$, then the corresponding value $y=f(x)$ is within a certain tolerance level of $L$.}
\begin{defn}[$\epsilon-\delta$ Definition of the Limit]
	Let $f$ be a function defined on all $x$ on the open interval $I$ containing $a$, except possibly at $a$ itself. The limit of $f(x)$ as $x$ approaches $a$ is $L$, defined as
	\begin{equation}
		\lim_{x\to a} f(x) = L.
	\end{equation}
	if for every $\epsilon > 0$, however small, there exists a $\delta > 0$ such that $|f(x)-L| < \epsilon$ whenever $0 < |x-a| < \delta$.
\end{defn}
As with a lot of mathematical concepts, it will help if we visualize these ``tolerances'' with our limit definition.

While epsilon-Delta proofs involve choosing a $\delta$ value first then conclude that $|f(x)-L|<\epsilon$, we actually have to do some \textit{scratch work}.

First, we begin with $|f(x)-L|<\epsilon$, then manipulate that expression algebraically to get something similar to $|x-a|<\ $\textit{something}, then choosing a $\delta$ value, which is what we begin the formal proof with. Let us prove Example \ref{eqn:firstexmp} rigorously.
\begin{exmp}
	Prove that $\displaystyle\lim_{x\to 2} 2x+1 = 5$.
	\begin{soln}
		Given the definition, we must show that $|2x + 1 - 5| < \epsilon$ whenever $0 < |x-2| < \delta$. To begin, we must first find the $\delta$ we need. \textit{Note that this is not part of the proper proof!}
		\begin{align*}
			|2x + 1 - 5|<\epsilon &= |2x - 4|<\epsilon \\
			&=2|x - 2|<\epsilon \\
			&=|x - 2|<\frac{\epsilon}{2}
		\end{align*}
		Choose $\delta=\epsilon/2$. With this $\delta$, we subsitute it back into the algebraic steps to prove the limit.
	\end{soln}
	\begin{proof}
		Given $\epsilon > 0$, choose $\delta = \dfrac{\epsilon}{2}.$ Assume that $0< |x-2| <\delta$. Then,
		\[ |(2x + 1) - 5| = |2x - 4| = 2|x - 2| < 2\delta = 2\left(\frac{\epsilon}{2}\right) = \epsilon\]
		Thus,
		\[\text{if }0 < |x-2| < \delta \text{ then }|(2x + 1) - 5| < \epsilon.\] Therefore, by definition of a limit, $\displaystyle\lim_{x\to 2} 2x+1 = 5$.
	\end{proof}
\end{exmp}

\section{Limit Theorems}
\begin{thm}\label{thm:lims}
	Given that the limits $\lim_{x\to a} f(x)$ and $\lim_{x\to a} g(x)$ exist, then\vspace{0.5em}
	\begin{enumerate}[itemsep=0.5em]
		\item For a constant $c$, $\displaystyle\lim_{x\to a}c=c$.
		\item For a real number $a$, $\displaystyle\lim_{x\to a}x=a$.
		\item $\displaystyle\lim_{x\to a}[f(x) + g(x)]=\lim_{x\to a} f(x) + \lim_{x\to a} g(x)$.
		\item $\displaystyle\lim_{x\to a}[f(x) - g(x)]=\lim_{x\to a} f(x) - \lim_{x\to a} g(x)$.
		\item $\displaystyle\lim_{x\to a}[c\,f(x)]=c \lim_{x\to a} f(x)$.
		\item $\displaystyle\lim_{x\to a}[f(x)\,g(x)]=f(x)\cdot g(x)$.
		\item If $\lim_{x\to a} g(x)\neq 0$, $\displaystyle\lim_{x\to a}\frac{f(x)}{g(x)}=\frac{\lim_{x\to a} f(x)}{\lim_{x\to a} g(x)}$.
		\item For a positive integer $n$, $\displaystyle\lim_{x\to a}[f(x)]^n= \left[\lim_{x\to a} f(x)\right]^n$.
		\item $\displaystyle\lim_{x\to a}\left[\sqrt[n]{f(x)}\right]=\sqrt[n]{\lim_{x\to a} f(x)}$ provided that the limit $L \in \mathbb{R}$.
	\end{enumerate}
\end{thm}
Fortunately, these theorems are intuitive. It is not unreasonable to assume that a function going to one limit, when added to another function with its own limit, results in the sum of their respective limits.
All of these theorems are provable using the Epsilon-Delta definition. In many other texts, these are called the \textit{Limit Laws} instead. 
\begin{exmp}\label{eqn:firstexmp}
	Evaluate $\displaystyle\lim_{x\to 2} 2x + 1$.
	\begin{soln}
		In this case, simply using the theorems will work.
		\begin{align*}
			\lim_{x\to 2} 2x + 1 &= 2\left(\lim_{x\to 2} x\right) + 1\\
			&= 2(2) + 1 \\
			&= 5\qedhere
		\end{align*}
	\end{soln}
\end{exmp}
To solve it in a more succint manner, we may directly substitute $x=2$ straight into $2x+1$.

In many cases, direct substitution of the value will lead to getting the correct answer of the limit. Evaluating manually by the table used earlier will verify this to be true.

\begin{exmp}
	Evaluate $\displaystyle\lim_{t\to \pi}\sin t$.
	\begin{soln}
		Similarly,
		\begin{align*}
			\lim_{t\to \pi}\sin t &= \sin\left(\lim_{t\to \pi} t\right)\\
			&= \sin \pi \\
			&= 0\qedhere
		\end{align*}
	\end{soln}
\end{exmp}

\begin{marginfigure}
	\centering
	\resizebox{\linewidth}{!}{%
		\begin{tikzpicture}[>=stealth, scale=0.8]
	% Axes
	\begin{axis}[
		axis lines = middle,
		xlabel = {$x$},
		ymin = -1.2, ymax = 1.2,
		xmin = 0, xmax = 6.5,
		domain = 0:2*pi,
		samples = 100,
		trig format plots=rad, 
		y axis line style={draw=none},
		xtick={3.1416},
		xticklabels={$\pi$},
		ytick=\empty,
		xticklabel style={yshift=-2mm}
	]
		
	\addplot[line width=1.2pt, mathPurple] {sin(x)};
	\addplot[dashed, gray!70] coordinates {(3.1416, -1.2) (3.1416, 1.2)};
	\end{axis}
\end{tikzpicture}%
	}
	\caption{$\sin(x)$ at $[0, 2\pi]$}
	\label{fig:sin-graph}
\end{marginfigure}

We see the function $\sin(x)$ graphed in Figure \ref{fig:sin-graph}. Interestingly, given the fact that the period of the sine wave is $2\pi$, its limit at each multiple of $\pi$ ($x$ values of $0,\ \pi,\ 2\pi,\ 3\pi,\ \ldots$) will always be 0.

Let us recall the first example in this chapter. Going back to the graphed function earlier,
\begin{exmp}
	Evaluate $\displaystyle\lim_{x\to 2} \dfrac{x^2-4}{x-2}$.
	\begin{soln}
		Direct substitution will result in $0/0$. This is undefined, for many reasons beyond the scope of this chapter.
		
		For functions where direct substitution fails, we need to find a way to simplify. Notice we can factor out $x^2-4=(x+2)(x-2)$. We may then directly substitute after. We should get 4 as the answer.
		\begin{align*}
			\lim_{x\to 2} \frac{x^2-4}{x-2} &= \lim_{x\to 2} \frac{(x+2)(x-2)}{x-2} \\
			&= \lim_{x\to 2} (x+2) \\
			&= \left(\lim_{x\to 2} x\right) + 2 \\
			&= 2+2 \\
			&= 4
		\end{align*}
		As is seen, the expected answer of 4 is obtained.
	\end{soln}
\end{exmp}

More ways to simplify will prove helpful.

\begin{exmp}
	Evaluate $\displaystyle \lim_{x\to 0} \left[\frac{2x^2+4x}{x}+\frac{x\sqrt{4+x}}{2}\right]$
	\begin{soln}
		By Theorem \ref{thm:lims}(3), split the function.
		\begin{align*}
			\lim_{x\to 0} \left[\frac{2x^2+4x}{x}+\frac{x\sqrt{4+x}}{2}\right] &= \lim_{x\to 0} \frac{2x^2+4x}{x} + \lim_{x\to 0} \frac{x\sqrt{4+x}}{2} \\
			&= \lim_{x\to 0} \frac{x(2x+4)}{x} + \lim_{x\to 0} \frac{x\sqrt{4+x}}{2} \\
			&= \lim_{x\to 0} (2x+4) + \lim_{x\to 0} \frac{x\sqrt{4+x}}{2} \\
			&= 2\left(\lim_{x\to 0} x\right) + 4 + \frac{\left(\displaystyle\lim_{x\to 0} x\right)\sqrt{4+\left(\displaystyle\lim_{x\to 0} x\right)}}{2} \\
			&= 2(0) + 4 + \frac{0\sqrt{4+0}}{2} \\
			&= 4\qedhere
		\end{align*}
	\end{soln}
\end{exmp}

It is assumed that the theorems are clear by this point. Showing the limit theorem application step-by-step for each solution will cause unnecessary length. As such, we will show shorter solutions. However, it is still good practice to properly apply the theorems in writing.

\begin{exmp}
	Evaluate $\displaystyle\lim_{t\to e} \frac{3\sqrt{h}}{h^4+\sqrt{3h}}$.
	\begin{soln}
		This is your reminder to look at the variables used in the limit, and what is in the function. The limit is as $t$ approaches $e$. However, there is no $t$ variable in the equation. By Theorem \ref{thm:lims}(1),
		\[\lim_{t\to e} \frac{3\sqrt{h}}{h^4+\sqrt{3h}}=\frac{3\sqrt{h}}{h^4+\sqrt{3h}}.\qedhere\]
	\end{soln}
\end{exmp}

\begin{exmp}
	Evaluate $\displaystyle\lim_{x\to -1}\frac{x^3+1}{x+1}$.
	\begin{soln}
		Direct substitution fails, since $\left((-1)^3+1\right)/\left(-1+1\right)=0/0$.
		
		Thus, we have to simplify by recalling factorization.\bmarginnote{Recall the sum of cubes which is $x^3 + a^3 = (x + a)(x^2 - ax + a^2)$.}
		\begin{align*}
			\lim_{x\to -1}\frac{x^3+1}{x+1} &= \lim_{x\to -1}\frac{(x+1)(x^2-x+1)}{x+1} \\
			&= \lim_{x\to -1} \left(x^2-x+1 \right) \\
			&= -1^2 - 1 + 1 \\
			&= 1\qedhere
		\end{align*}
	\end{soln}
\end{exmp}


\section{One-Sided Limits}
The \textit{Heaviside function} $H$, named after Oliver Heaviside (1850-1925), represents a signal that switches on at a specified time and stays switched on indefinitely.

\begin{equation}
	H(x) := \begin{cases}
		1 & x \geq 0 \\
		0 & x < 0
	\end{cases}
\end{equation}

\begin{marginfigure}
	\centering
	\resizebox{\linewidth}{!}{%
		\begin{tikzpicture}[>=stealth, scale=0.8]
	% Axes
	\draw[->] (-4, 0) -- (4, 0) node[below] {$x$};
	\draw[->] (0, -1) -- (0, 3) node[left] {$y$};
	
	\draw (-0.3, 1.25) node[below, font=\footnotesize] {$1$};
	\draw (-0.2, 0) node[below, font=\footnotesize] {$0$};
	
	\draw[line width=1.2pt, mathPurple] (-4, 0) -- (-0.02, 0);
	\draw[line width=1.2pt, mathPurple] (0, 1) -- (4, 1);
	\draw[line width=1.2pt, mathPurple, fill=white] (0,0) circle (2.5pt);
	\draw[line width=1.2pt, mathPurple, fill=mathPurple] (0,1) circle (2.5pt);
\end{tikzpicture}%
	}
	\caption{The \textit{Heaviside function} $H(x)$, where it breaks at $x=0$.}
	\label{fig:heaviside-func}
\end{marginfigure}

It is interesting because if we take the limit of $H(x)$ at $x=0$, what do we get? Remember, the value at 0 is not automatically its limit at 0.

From the left (negative values), we see the function approaches 0. This is denoted as
\[\lim_{x\to 0^-} H(x) = 0.\]
From the right (positive values), the function approaches 1. This is denoted as
\[\lim_{x\to 0^+} H(x) = 1.\]

The previous principles and theorems discussed still hold for one-sided limits. In relating the ordinary limit and one-sided limits,

\begin{thm}
	For a function $f(x)$, its limit at point $a$, $\lim_{x\to a} f(x)$ exists if and only if both of its one-sided limits $\lim_{x\to a^-} f(x)$ and $\lim_{x\to a^+} f(x)$ also exist.
\end{thm}




\section{Continuity of a Function}
\begin{defn}[Continuity at $a$]\label{defn:continuity}
	A function $f(x)$ is said to be continuous at $a$ when $f(a)$ exists, $\lim_{x\to a} f(x)$ exists, and $\lim_{x\to a} f(x) = f(a)$.
\end{defn}
Further, we may say that if $f(x)$ is continuous on $(-\infty, \infty)$, the function is continuous everywhere.
\begin{thm}
	For two continuous functions $f$ and $g$, The functions $f+g$, $f-g$, $fg$, and $f/g$ (only when $g\neq 0$) are also continuous.
\end{thm}
\begin{proof}
	To begin this proof, we will rely on the aforementioned limit theorems. Let us focus on the addition of the two functions.
	
	By Definition \ref{defn:continuity}, we need to show that
	\[\lim_{x\to a}[(f+g)(x)] = (f+g)(a).\]
	
	Using the definition of function addition, $(f+g)(x) = f(x) + g(x)$. Thus,
	\begin{align*}
		\lim_{x\to a}[f(x)+g(x)] &= (f+g)(a) \\
		\lim_{x\to a} f(x) + \lim_{x\to a} g(x) &= (f+g)(a) \\
		f(a) + g(a) &= (f+g)(a) \\
		(f + g)(a) &= (f+g)(a)
	\end{align*}
	The proofs for the other three function operations are left as exercises to the reader.
\end{proof}